\paragraph{Structure of the \(C_2\) minors.}
Write a coordinate row of \(C_2\) as
\[
  (1,y_1,y_2,y_1y_2),
\]
where \(y_1\) is the structural point chosen at the first gate and \(y_2\) is
the point chosen at the child gate.  For \(k=0\), \(L_2=D_2=4\), and the
unique maximal minor factors as
\[
  \det C_2
  = \pm (t_{\varnothing}^{1}-t_{\varnothing}^{0})^2
      (t_{0}^{1}-t_{0}^{0})
      (t_{1}^{1}-t_{1}^{0}).
\]
Thus sibling distinctness is sufficient in the square case.

For \(k=1\), \(L_2=8\) and \(D_2=4\), so there are
\(\binom{8}{4}=70\) maximal minors.  Exhaustive symbolic factorization
splits them into three classes according to the multiplicities with which
the four lower blocks are selected.  The \(6\) minors of type \((2,2)\)
factor into two linear point differences and one squared linear
difference.  The \(48\) minors of type \((2,1,1)\) factor into four linear
point differences.  The remaining \(16\) minors, of type
\((1,1,1,1)\), contain one coordinate from each lower block and give a
single degree-\(4\) relation involving \(8\) structural points; the
representative polynomial is irreducible over \(\mathbb{Q}\).
A numerical sanity check over \(\mathbb{F}_{2^8}\) tests every candidate
point-difference factor \(u-v\) by imposing \(u=v\) and randomly
specializing the remaining variables.  Nonzero specializations rule out
such divisibility in characteristic two.  Consequently, pairwise
distinctness of structural points explains all minors containing a double
block but does not account for the hard four-single class.  Since this
class has \(2^{D_i}\) child choices in the analogous higher-level
configuration, the direct minor-certificate route does not presently
yield a polynomial-size checked setup for \(C_4\) or \(C_5\).
